\chapter{Landasan Kepustakaan}
\section{Kajian Pustaka}
Bandingkan penelitian terdahulu berdasarkan masalah, data, metode, evaluasi,
hasil, dan keterbatasannya. Sintesis harus menunjukkan posisi penelitian ini.
Gunakan sitasi nama--tahun dengan koma, misalnya \parencite{lamport1994latex}.

Langkah sintesis dapat diringkas sebagai berikut:
\begin{enumerate}
  \item kelompokkan penelitian berdasarkan pertanyaan riset;
  \item bandingkan data, metode, dan evaluasi; dan
  \item tarik kesenjangan yang mendasari penelitian yang diusulkan.
\end{enumerate}

\section{Landasan Teori}
Jelaskan konsep yang benar-benar digunakan dalam metodologi. Persamaan contoh:
\begin{equation}
  \bar{x}=\frac{1}{n}\sum_{i=1}^{n}x_i.
  \label{eq:rata-rata}
\end{equation}
Persamaan~\ref{eq:rata-rata} menunjukkan cara menghitung rerata sampel.

\begin{figure}[htbp]
  \centering
  \fbox{\parbox[c][2.4cm][c]{0.78\textwidth}{\centering
    Masukan penelitian $\longrightarrow$ Metode $\longrightarrow$ Evaluasi}}
  \caption{Contoh diagram alur konseptual penelitian}
  \label{fig:alur-konseptual}
\end{figure}
